\section{Introduction}
A finite machine instruction reads a short part of a configuration and
changes it. A material flow, by contrast, must be invertible at every
finite time. Two questions therefore arise when a particle is to carry a
computation: where is the information erased by an instruction retained,
and how is the resulting injective update made compatible with
incompressibility? We study these questions for finite prefix instructions
and their smooth three-dimensional realizations.

A stack is an infinite word whose first letter is its accessible end.
A prefix instruction has the form
\[
 (s;v,w)\longmapsto(s';v',w'),
\]
meaning that it sends $(s,vL,wR)$ to $(s',v'L,w'R)$ for every pair of
unchanged tails $L,R$. Here $s,s'$ belong to a finite set of control
states, and all four prefixes are finite words over a finite alphabet.
A positional code turns this operation into a positive diagonal affine
map between two rectangles. The source rectangles must be separately
disjoint, and so must the target rectangles. Intersections between a
source rectangle and a target rectangle are allowed. These conditions
are stronger than injectivity along one initialized computation.

The planar map need not preserve area. Appending one history letter,
for example, contracts one stack coordinate without expanding the
other. We realize such instructions in two ways. A compensating normal
strain transports the entire coded rectangle and a three-dimensional
neighborhood. Alternatively, a compressible planar motion is placed in
an invariant plane of an incompressible field. The latter preserves the
prescribed planar dynamics without prescribing an affine map off the
plane. Both methods require control of the complete continuous path:
correctness at integer sampling times alone does not establish a fixed
particle observation.

\subsection{The fluid problem and a first result}
Let $\mathbb T^3=(\mathbb R/\mathbb Z)^3$ with its flat metric, and fix a
positive computable viscosity $\nu$. For an arbitrary fixed positive real
$\nu$, the same formulas and conclusions hold, and numerical evaluation
is relative to that coefficient. We use
\begin{equation}\label{eq:sheets-ns}
 u_t+(u\cdot\nabla)u-\nu\Delta u+\nabla p=f,
 \qquad \operatorname{div}u=0,\qquad u(0)=0.
\end{equation}
An instance is a deterministic machine with finite state set, finite work
alphabet with a blank, an initial state, a specified halt set, and a
finite input written on an otherwise blank two-sided tape. Its head
starts at cell zero. Missing transitions may be interpreted as halting.
For a fixed initial label $a$, the material position satisfies
$X'(t)=u(t,X(t))$ and $X(0)=a$. The observation asks whether
$X(t)\in O$ for some $t\ge0$, where $O$ is a fixed open set.

A finite effective force prescription is a finite program that evaluates
$f$ and any requested mixed derivative at computably specified arguments
to any positive rational error, and supplies the derivative bounds used
on compact sets. The program describes every instruction branch. It does
not generate a machine execution and replay that execution as a force.
On a torus our classical comparison class requires $u,u_t$, spatial
derivatives of $u$ through order two, and $p,\nabla p$ to be continuous
on every finite closed time cylinder. Velocity and pressure are periodic,
and pressure has spatial mean zero. The prescribed solutions themselves
are smooth.

\begin{theorem}[Interspersed history on moving sheets]\label{thm:sa-interleaved}
At fixed positive computable viscosity, every machine and finite input
effectively determine a smooth mean-zero general force on the unit
$\mathbb T^3$, with all mixed derivatives bounded and period one for
$t\ge1$.  Its zero-data solution is globally smooth and unique in the
classical periodic comparison class, with uniformly bounded kinetic
energy.  Its material trajectory from $(1/8,1/4,1/4)$ enters
$\{x:1/2<x_1<1\}$ at some finite time exactly when the machine halts.
The repeated template depends only on the table, and the input enters
only the loading pulse.
\end{theorem}


This first result is proved with history records interspersed among work
letters. Its geometric map changes planar area, so a third scale restores
volume. Later constructions change the history arrangement or initialization
when a different property is wanted: compact Euclidean support, a bounded
detector, periodicity from the initial time, or a spatial clock. The
comparison classes, initial labels, detectors and temporal conclusions are
stated separately for these alternatives.

\subsection{Antecedents and the role of the present constructions}
The logical obstruction comes from Turing's undecidable symbol-printing
problem~\cite[Section~8]{Turing1936}. Modify such a machine to halt at
the first printing of the designated symbol, and redirect any earlier
ordinary stop into an infinite dummy loop. It then halts exactly when
the designated symbol would have been printed, giving the halting
formulation used here. Retaining information lost by an irreversible instruction was
discussed by Landauer~\cite[Section~3]{Landauer1961} and implemented by
Bennett's history recording~\cite[Eqs.~(9)--(11)]{Bennett1973}. Our
histories remain stored. We prove the local tables and their inverses
explicitly; the cleanup and complexity statements of Bennett's reversible
simulation are not used.

Moore's generalized shifts connect finite symbolic instructions with
positional codes and piecewise affine dynamics
\cite{Moore1990,Moore1991}. The smooth realization and suspension for
invertible generalized shifts in
\cite[Section~6, Theorem~12 and its corollary]{Moore1991} are direct
antecedents of the present viewpoint. We need additional statements on
full closed rectangles, effective derivative bounds, incompressibility,
zero initial velocity, and observation at every intermediate real time.
These are established by the explicit motions below. The classical
fragility of positional tape storage also explains why an exact event
should not be read as a uniform finite-precision tolerance.

Cardona, Miranda, Peralta-Salas and Presas constructed Turing-complete
stationary Euler particle flows on a three-sphere with a chosen metric
\cite{CardonaMirandaPeraltaSalasPresas2021}. Their area-preserving
extension of generalized shifts uses contraction, transport and expansion
of separated blocks followed by a relative area correction
\cite[Proposition~5.1]{CardonaMirandaPeraltaSalasPresas2021}.
That construction is an important predecessor of the planar routing used
here. Our explicit normal compensation and invariant-plane formulas give
the prescribed flat three-dimensional motions directly; their full
neighborhood and detector properties are proved at the point of use.
Cardona, Miranda and Peralta-Salas also constructed Turing-complete
stationary Euler flows for the standard Euclidean metric on $\R^3$,
with computation on a noncompact invariant plane and infinite kinetic
energy \cite[Theorem~1 and the discussion following it]{CardonaMirandaPeraltaSalas2023Beltrami}.

Universality is also known with viscosity. Dyhr, Gonz\'alez-Prieto, Miranda
and Peralta-Salas construct unforced stationary Navier--Stokes particle
universality on suitable compact three-manifolds after deforming the metric,
using Hodge viscosity and a harmonic velocity
\cite[Theorem~A]{DyhrGonzalezPrietoMirandaPeraltaSalas2026}.
The fixed flat domains, prescribed external force, and zero initial
velocity in the present paper specify a different fluid experiment.
Our results concern these constructed solutions and exact material events;
they make no assertion about arbitrary-data regularity.

The companion \emph{Finite Instructions and Solenoidal Shear Flows}
\cite{OpenAIShear2026} gives a complete initialized construction with
reciprocal planar maps and directly solenoidal force at pressure zero.
We use its full-domain recorder in one later specialization, with an
explicit identification of all guards. The general residual forces below
need not be solenoidal. When a torus force is projected, the accompanying
pressure changes; that operation is kept distinct from the direct shear
identity.

\subsection{Proof strategy and reading order}
The first application has four steps. Gapped radix intervals turn complete
prefix cylinders into separated closed rectangles. A local curl formula
then implements an affine volume-preserving motion on a neighborhood.
Private heights separate the instruction rectangles while a reciprocal
normal strain corrects their planar area changes. Finally, a history
invariant identifies the computation, a loading segment places the fixed
particle at its initial code, and the horizontal interpolation proves
the detector equivalence throughout every instruction period.

Section~\ref{sec:model} records the residual realization and the exact
analytic comparison classes. Section~\ref{sec:sa-sheets} develops the
radix and moving-sheet tools. Section~\ref{sec:sa-histories} completes
Theorem~\ref{thm:sa-interleaved}. Sections~\ref{sec:sheets-boxes}
and~\ref{sec:sheets-planes} then explain the additional requirements of
full boxes and invariant-plane routing. Sections~\ref{sec:sa-separator-histories}--\ref{sec:sa-marked-five-stage}
organize the resulting processors by where they retain history: next to
a separator, on a tape track, or in a separate stack.
Sections~\ref{sec:sb-prefix}--\ref{sec:sb-events} treat further
history placements, boot instructions and fixed observers.
Section~\ref{subsec:p2:rapid-periodic} gives two compact reciprocal
realizations. The final section replaces a temporal program by a spatial
clock and verifies its separate decay and phase-window observation.
